\begin{lemma}[Change lemma, mixed location: final starred stage closed]
Fix a finite loopless multigraph with its exact edge set, a palette of
$k$ colors, and a distinguished edge $e_0$.  Let $(T,\varphi_n)$ be a
positive-rung extended Tashkinov tree.  Let $\mathcal B$ be a hierarchy
base and let
\[
 T_n=T_{n,0}\subset T_{n,1}\subset\cdots\subset
 T_{n,q}\subset T_{n,q+1}=T
\]
be a good hierarchy under $\varphi_n$, where $q$ is its last divider
index and
\[
 A=T_{n,q}^{*}=\noderef{HierarchyStarredStage}(H,q).
\]
Assume that $T$ has maximum property, that $V(T)$ is not elementary,
that $\varphi_n$ is proper and leaves exactly $e_0$ uncolored, and that
the rung number is positive.  These data form a
\noderef{HierarchyCounterexample} $C$.

Assume the six lower-rung assertions of the technical theorem, uniformly
for every positive-rung maximum ETT on the same graph, palette, and
distinguished edge, with arbitrary certified input and generating
colorings and with rung number smaller than that of $T$: lower-rung
elementarity, terminal interchangeability, the two PE stability and
stored-pair assertions, the terminal PE component assertion, the
defective-endpoint order assertion, and the mod-coloring assertion.
Assume also the same-rung induction hypothesis that every maximum ETT
with a good hierarchy having fewer than $q$ dividers is elementary.
Finally assume the large-color estimate
\[
 2n+11\leq |\overline\varphi_n(T_n)|.
\]

Put $\ell=|T|-|A|$, the number of exact tail additions above $A$, and let
\[
 T(y_{\ell-1})=T\big|_{\min\{|A|+(\ell-1),|T|\}}
\]
be the exact ordered prefix supplied by \noderef{TreeSequencePrefix}.
Assume (6.6): $V(T(y_{\ell-1}))$ is elementary under $\varphi_n$.
Assume moreover that $C$ is lexicographically minimal among all
\noderef{HierarchyCounterexample}s on the fixed graph, palette, and
$e_0$ which have exactly $q$ dividers, allowing their certified input
colorings, generating colorings, bases, ETTs, and good hierarchies to
vary.  Explicitly, no such counterexample $C'$ satisfies
\[
 \noderef{HierarchyCounterexampleLess}
  \bigl(\noderef{HierarchyCounterexampleMeasure}(C'),
        \noderef{HierarchyCounterexampleMeasure}(C)\bigr).
\]

Let $\alpha\ne\beta$ be colors missing on $T(y_{\ell-1})$.  Write
\[
 i_c=\noderef{HierarchyFirstMissingPosition}(\varphi_n,T,c),
 \qquad v_c=T.\mathit{vertex}(i_c).
\]
If at least one of $\alpha,\beta$ is not missing on $A$, assume both
\[
 v_\alpha\prec v_\beta                                      \tag{a}
\]
and the exact edge-identity exclusion
\[
 \alpha\notin
 \left\{c:\exists t\in\operatorname{Fin}|T|,
       |A|\leq t<i_\beta,
       \ \varphi_n(T.\mathit{edge}(t))=c\right\}.             \tag{b}
\]
Thus (b) is precisely
$\alpha\notin\varphi_n\langle T(v_\beta)-A\rangle$; it is not
an endpoint-only abbreviation.


Assume in addition that the first $\alpha$-missing vertex lies in the
last starred stage and the first $\beta$-missing vertex lies outside it:
\[
        v_\alpha\in V(A),\qquad v_\beta\notin V(A).
\]
Assume that $\alpha$ is missing on $A$ and that the singleton
$\{\alpha\}$ is closed on $A$ under $\varphi_n$.  Then
\[
 \noderef{TwoColorComponent}(G,\varphi_n,\alpha,\beta,
       \operatorname{inl}v_\alpha)
 =
 \noderef{TwoColorComponent}(G,\varphi_n,\alpha,\beta,
       \operatorname{inl}v_\beta).
\]
Moreover, $\alpha$ and $\beta$ are $A$-interchangeable: among all paper
$(\alpha,\beta)$-path components, including isolated vertex components,
at most one meets $A$.
\end{lemma}
\begin{proof}
Write $H=C.\mathit{hierarchy}$, put
$q=H.\mathit{dividerCount}$, and retain the notation $A$, $\ell$,
$T(y_{\ell-1})$, $i_\alpha,i_\beta,v_\alpha,v_\beta$ from the
statement.  Write
\[
 K_x=\noderef{TwoColorComponent}
       (G,\varphi_n,\alpha,\beta,\operatorname{inl}x).
\]
Thus $K_x$ is the full component in the incidence graph and contains the
individual edge nodes of the multigraph.

We first record the prefix facts used below.  The nestedness and terminal
fields of \noderef{TashkinovHierarchy}, together with the two cases in
\noderef{HierarchyStarredStage}, make $A$ an exact ordered initial
segment of $T$.  If $\ell=0$, then $|A|=|T|$ and
\noderef{TreeSequenceInitialSegment} identifies $A$ with $T$ entry by
entry, including every exact edge identity.  The prefix
$T(y_{\ell-1})$ from \noderef{TreeSequencePrefix} would then also be
$T$.  Its assumed elementarity would contradict the
\texttt{notElementary} field of \noderef{HierarchyCounterexample}.  Hence
$\ell\geq1$.  It follows that $A$, and every ordinary or starred
hierarchy stage at or below $q$, is an ordered initial segment of
$T(y_{\ell-1})$.  Restricting
\noderef{ColoredVertexSetElementary} makes each of these stages
elementary.

The least-candidate definition
\noderef{HierarchyFirstMissingPosition}, applied to the two colors
missing on $T(y_{\ell-1})$, says that $v_\alpha$ misses $\alpha$ and
$v_\beta$ misses $\beta$.  Their stated locations make them distinct.
Elementarity of the prefix then says that $v_\alpha$ does not miss
$\beta$ and $v_\beta$ does not miss $\alpha$.  By
\noderef{MissingColors} and \noderef{PresentColors}, each is incident
with an edge of the other color.  Properness therefore makes $K_{v_\alpha}$
and $K_{v_\beta}$ nontrivial path components in the sense of
\noderef{TwoColorNontrivialPathComponent}; this also explains why the
singleton-inclusive convention in \noderef{TreeInterchangeable} causes
no endpoint exception here.

We next prove interchangeability on $A$.  If $q>0$, then
$A=T_{n,q}$ by \noderef{HierarchyStarredStage}.  The last-stage instance
of \noderef{step1} applies: the statement supplies maximum property,
the six lower-rung assertions, the fewer-divider hypothesis, the
large-color bound, elementarity of the full-tail prefix, and the
missing and singleton-closed color $\alpha$.  If $q=0$ and the terminal
extension is PE, conclusion (iii) of \noderef{interchange} applies to
the exact PE core $A$; the identity recoloring satisfies both clauses
of \noderef{CombinedStableColoring}.  If $q=0$ and the terminal
extension is SE or RE, \noderef{HierarchyBase} identifies $A$ with
$T_n$.  Restricting the certified series to the preceding rung and its
maximum-property data allows the lower-rung terminal-interchangeability
assertion to be used when the rung number exceeds one; the SE stable
relation is reversed by \noderef{sc2}, while in RE the generating
coloring is unchanged.  At rung one the same assertion is the
interchangeability conclusion of \noderef{TechnicalTheoremBase}.
Thus
\[
                 \alpha,\beta\text{ are }A\text{-interchangeable}.
                                                               \tag{1}
\]

Suppose for contradiction that $K_{v_\alpha}\ne K_{v_\beta}$.
The component $K_{v_\alpha}$ meets $A$ at $v_\alpha$.  By (1),
\noderef{TreeInterchangeable}, and
\noderef{TwoColorComponentsDisjoint}, the component $K_{v_\beta}$ is
vertex-disjoint from $A$ and has no $\alpha$- or $\beta$-colored edge
incident with $A$.  Let
\[
       \sigma_n=\varphi_n/K_{v_\beta}
\]
be the full component switch of \noderef{TwoColorComponentSwap}.
The theorem \noderef{TwoColorComponentSwapProper} makes $\sigma_n$
proper and the switch does not alter the unique uncolored edge, so
\noderef{IsOneEdgeUncolored} still holds.  It leaves
$v_\alpha$ missing $\alpha$ and makes $v_\beta$ missing $\alpha$.

Let $B=T(v_\beta)$ be the exact ordered prefix of $T$ ending at the
entry $v_\beta$.  Since $v_\beta$ lies in $T(y_{\ell-1})-A$, if
$s=|B|-|A|$, then
\[
                         1\leq s\leq\ell-1<\ell.               \tag{2}
\]
Both $v_\alpha$ and $v_\beta$ lie in $B$, so $B$ is not elementary
under $\sigma_n$.

Before transporting the coloring, form the old-color ETT on this exact
prefix.  Restrict the final certified
\noderef{TashkinovClosure} of the given ETT at the displayed length
$|B|$, retaining the entire ladder series and its exact edges.  The
terminal ladder tree is an initial segment of $A$, while
$A\prec B\prec T$, so the terminal containment fields remain strict.
This gives an ETT $\mathcal E_B$ on the actual ordered tree $B$ under
$\varphi_n$.  Its rung and ladder data are unchanged; consequently its
\noderef{MaximumProperty} witness is the one carried by $C$.  Restrict
the terminal hierarchy stage to $B$ and retain the same base, divider
positions, starred stages, and Gamma pairs.  Every such stage lies in
$A$.

The switch is outside $A$.  Applying
\noderef{OutsideKempeRecoloringStable} first to the terminal ladder tree
and, in PE, separately to both constituents of the exact core, gives
terminal stability and the required
\noderef{CombinedStableColoring}.  In PE,
\noderef{CombinedStableClosurePreserved} transports the same closure
$R_n$ and its individual augmentation edges.

We check the route from the core to $B$ before invoking a correspondence
theorem.  The ordered-unused hypothesis gives
\[
 \alpha\notin\varphi_n\langle B-A\rangle,                     \tag{3}
\]
and the definition of $i_\beta$ gives
$\beta\notin\overline\varphi_n(B-v_\beta)$.  Each ordinary displayed
edge of $B-A$ is supplied by \noderef{TashkinovAugmentation}; if it had
color $\beta$, its preceding prefix would already contain a vertex
missing $\beta$, contrary to the minimality of $i_\beta$.  The only
nominal exception is the exact first SE/RE connecting seed when
$q=0$, $A=T_n$, and $\beta\in D_n$.  That individual seed is incident
with $T_n$ by the SE or RE constructor of
\noderef{TashkinovSeriesStep}, so the incident-edge clause of the
terminal \noderef{StableColoring} fixes its color.  Thus (3), the
first-missing property, and this exact first-seed check show that every
displayed edge used from the appropriate PE or SE/RE seed to $B$ has
the same color under $\sigma_n$ and $\varphi_n$.  At every old
augmentation witness before $v_\beta$, the switched component either
does not meet the vertex or exchanges two present incident colors; the
missing set is therefore unchanged.  The exact old augmentation route
to $B$ remains valid.  Equivalently, its required restriction is the
one supplied by
\noderef{TashkinovAugmentationSequenceRestrict}.

Now apply \noderef{Td} to the terminal stable relation, using the given
six lower-rung assertions, and obtain the required
\noderef{ModColoring}.  In PE, apply \noderef{HierarchySplitter} to
the old-color prefix ETT $\mathcal E_B$, using the combined stability,
transported closure, and verified exact augmentation sequence.  It
returns a \noderef{CorrespondingExtendedTashkinovTree}, maximum property,
and the same hierarchy stages.  In SE or RE, the exact one-edge seed
and its verified route are an eligible tree in the universal clause of
\noderef{ModColoring}; its corresponding ETT retains the non-PE kind,
and \noderef{MaximumPropertyTransport} transports maximum property.

It remains to construct the good hierarchy; it is not inferred from the
correspondence certificate.  Since every hierarchy stage lies in $A$
and the switched component is disjoint from $A$, every missing set at a
stage vertex and every exact edge incident with a stage is unchanged.
Use the old Gamma pairs.  Their two-element size and pairwise
disjointness are literal.  Avoidance holds because every segment edge
keeps its exact color.  The unique-refresh field holds because both the
stage differences and all their vertexwise missing sets are unchanged.
In PE, the two constituent stable relations and
\noderef{CombinedStableClosurePreserved} preserve both clauses of the
parallel-base field.  Finally, every boundary edge relevant to every
positive-stage minus-closed field is incident with its stage and hence
fixed.  The old minus-closed witnesses therefore prove the new ones.
These five verifications construct a
\noderef{GoodHierarchy} witness under $\sigma_n$.  Only after constructing
that witness do we compare it with the old witness by
\noderef{SameGoodHierarchy}; the stages, stored sets, and Gamma pairs
are pointwise identical.

The resulting ETT on $B$, its maximum property, base, hierarchy,
newly constructed good witness, and the two vertices missing $\alpha$
form a \noderef{HierarchyCounterexample} $C'$ with the same divider
count as $C$.  We finally compare the actual measures.  Put
$r_0=\noderef{HierarchyPathNumber}(H)$ and let $r'_0$ be the path
number of the hierarchy truncated at $B$.  If $s<r_0$, the fallback
clause gives $r'_0\leq s<r_0$.  If $r_0\leq s$, the path tail beginning
at $r_0$ restricts to a path tail in $B$, with a singleton allowed, so
$r'_0\leq r_0$.  In the equality case the second coordinate drops:
if $T_{n,q}$ is the last ordinary stage and
$d=|A|-|T_{n,q}|$, then
\[
 |T|-|T_{n,q}|=d+\ell,\qquad |B|-|T_{n,q}|=d+s<d+\ell.
\]
Hence
\[
 \noderef{HierarchyCounterexampleLess}
   \bigl(\noderef{HierarchyCounterexampleMeasure}(C'),
         \noderef{HierarchyCounterexampleMeasure}(C)\bigr)
\]
holds.  This contradicts the stated minimality of $C$.  Therefore
$K_{v_\alpha}=K_{v_\beta}$, and (1) supplies the asserted
singleton-inclusive interchangeability.
\end{proof}

